% !PRINTED 703
\weberpage{703}
%<<CONT
einen Beitrag zu dem Gliede mit $z_1^{-1}$, und dieser hat nach (2) den Koeffizienten $a_{-1}$. Daraus folgt:

2. \gesperrt{Das Residuum von $dJ$ ist von der Wahl der Funktion $z$ unabhängig.}

Es gilt ferner noch der folgende Satz:

3. \gesperrt{Die Summe aller Residuen eines Differentials ist gleich Null.}

Wir erweitern den Ausdruck $\mathfrak{A}/\mathfrak{B}$ des Differentials $dJ$, wenn nötig, durch Hinzufügung von Punkten im Zähler und Nenner, so daß die voneinander verschiedenen Punkte des Unterecks
\weq{3}{\mathfrak{P}_1 \mathfrak{P}_2 \ldots \mathfrak{P}_n}
ein Polygon einer eigentlichen Klasse bilden. Die Pole von $dJ$ sind dann unter diesen $\mathfrak{P}$ enthalten, möglicherweise auch noch andere Punkte, für die das Residuum von $dJ$ gleich Null ist.

Als Nenner von $dJ$ können wir dann
\weq{4}{\mathfrak{B} = \mathfrak{P}_1^{m_1} \mathfrak{P}_2^{m_2} \ldots \mathfrak{P}_n^{m_n}}
nehmen, wenn wir die zuviel genommenen Punkte wieder im Zähler zufügen.

Wir nehmen eine Funktion $z$ in $\Omega$ von der $n$ten Ordnung, die in jedem der Punkte $\mathfrak{P}_i$ und in jedem nur zur ersten Ordnung unendlich wird, und bilden die Entwickelung für den Punkt $\mathfrak{P}_i$:
\weq{5}{\frac{dJ}{dz} = a_{m-2}^{(i)}z^{m-2} + \cdots + a_0^{(i)} + a_{-1}^{(i)}z^{-1} + \eta^{(i)}z^{-2}.}

Wenn wir für den Koeffizienten $a_r^{(i)}$ den Wert 0 zulassen, können wir in allen Punkten $\mathfrak{P}_i$ mit derselben Potenz $z^{m-2}$ anfangen. Die Summe der Residuen von $dJ$ ist
\[
-\overset{i}{\sum} a_{-1}^{(i)}.
\]

Nach § 182, 11. läßt sich ein Funktionensystem $\alpha_1, \alpha_2, \ldots, \alpha_n$ so bestimmen, daß $\alpha_r$ in dem Punkte $\mathfrak{P}_r$ unendlich klein in der ersten Ordnung wird, in den übrigen Punkten $\mathfrak{P}_i$ endlich und von Null verschieden bleibt.

Setzt man dann
\[
\alpha_1^m \alpha_2^m \ldots \alpha_n^m = \alpha_i^m \varrho_i,
\]
so wird
\weq{6}{\begin{aligned}
\varrho_1 &= 0^m \text{ in } \mathfrak{P}_2, \mathfrak{P}_3, \ldots, \mathfrak{P}_n, & \text{endlich und nicht} &= 0 \text{ in } \mathfrak{P}_1, \\
\varrho_2 &= 0^m \text{ ,, } \mathfrak{P}_1, \mathfrak{P}_3, \ldots, \mathfrak{P}_n, & \text{,, ,, ,,} &= 0 \text{ ,, } \mathfrak{P}_2, \\
&\cdots\cdots\cdots\cdots\cdots\cdots\cdots\cdots \\
\varrho_n &= 0^m \text{ ,, } \mathfrak{P}_1, \mathfrak{P}_2, \ldots, \mathfrak{P}_{n-1} & \text{,, ,, ,,} &= 0 \text{ in } \mathfrak{P}_n.
\end{aligned}}
